\scanpage{94}
\mspage{11}
% archive sheet number: 93

\[
H^{\unsure{*|\alpha}}(\bar U,\Lambda_\nu)\Longleftarrow E_2^{pq}=H^p(\tilde X_0,\underline\Psi_\nu^q)
\tag*{(2.5)}
\]
avec
\[
\underline\Psi_\nu^q=R^qg_*((\Lambda_\nu)_{\bar U}\mid\tilde X_0)=\varinjlim_\alpha R^qg_{\alpha*}((\Lambda_\nu)_{U_\alpha})\mid\tilde X_0,
\tag*{(2.6)}
\]
où on a écrit $\bar K=\varinjlim K_\alpha$ ($K_\alpha$ extensions finies de $\bar K/\tilde K$), $g_\alpha:\deleted{\unsure{U_\alpha}}\tilde U\mathbin{\times}_{\tilde S}S_\alpha\longrightarrow\tilde X$ \deleted{\unsure{et}} l'immersion canonique.

\deleted{\unsure{2.7.}} 2.7. Ici le calcul des faisceaux $\underline\Psi_\nu^q$ est de nature locale sur $\tilde X$. Notons d'ailleurs \unsure{leurs|les} relations de ces faisceaux avec \deleted{\unsure{ceux de}} les faisceaux $\underline{\unsure{\Phi|\varphi}}_\nu^q$ du n° 1 (que nous n'utilisons pas)
\[
\left\{\begin{aligned}
&\underline\Phi_\nu^q=\underline\Psi_\nu^{q-1}\qquad\text{si}\quad q\geq2\\
&0\longrightarrow\underline\Phi_\nu^0\longrightarrow(\Lambda_\nu)_{X_0}\longrightarrow\underline\Psi_\nu^0\longrightarrow\underline\Phi_\nu^1\longrightarrow0\ .
\end{aligned}\right.
\tag*{(2.8)}
\]
Donc : \unsure{peu de chose près}, ces faisceaux $\underline\Psi_\nu^{\unsure{i|q}}$ sont les faisceaux de cycles évanescents pour $\tilde X/\tilde S$ et l'ouvert $\bar U$ de $\tilde X_{\tilde{\unsure{\eta|y}}}$ (coeff les $\Lambda_\nu$).

2.9. Il y a lieu de considérer (2.5) comme une suite spectrale de $\unsure{\pi|\Pi}$-modules, et l'expression (2.6) comme un \unsure{nom|hom|isom.} de faisceaux \inserted{sur $\tilde X_0$} où $\pi$ opère, (de façon compatible avec ses opérations sur $\tilde X_0$ via $\pi_0$). Nous nous intéressons d'abord à l'action du sous-groupe \deleted{$I$} d'inertie $I$.

2.10. On peut donner encore une variante évidente des développements précédents, en remplaçant $\bar K$ par $\bar K_t$, cf 1.26. On trouve
